\documentclass[aspectratio=169]{beamer}
\usetheme{Madrid}
\usecolortheme{default}

% Packages for math and graphics
\usepackage{amsmath}
\usepackage{graphicx}
\usepackage{hyperref}

% Title Information
\title[Trajectory Optimization]{Introduction to Trajectory Optimization}
\subtitle{From Analytical Maneuvers to Numerical Solutions}
\author{Astrodynamics Module}
\date{Week 1}

\begin{document}

% --- Title Slide ---
\begin{frame}
    \titlepage
\end{frame}

% --- Objective Slide ---
\begin{frame}{Module Objective}
    \textbf{The Shift from Analytical to Numerical}
    \vspace{0.5cm}
    \begin{itemize}
        \item Transition from basic analytical maneuvers (e.g., Hohmann transfers) to numerical optimization.
        \item Learn how modern mission designers calculate complex trajectories.
        \item \textbf{Goal:} Minimize fuel ($\Delta V$) or time while satisfying complex, non-linear mission constraints.
    \end{itemize}
\end{frame}

% --- Day 1 ---
\begin{frame}{Day 1: The Core Conflict}
    \framesubtitle{Targeters vs. Optimizers}
    
    \begin{columns}
        \begin{column}{0.5\textwidth}
            \textbf{Targeters (Root-Finding)}
            \begin{itemize}
                \item E.g., Differential Correctors.
                \item Attempt to drive a mathematical error \textit{exactly} to zero.
                \item If constraints are too tight, they diverge or fail.
                \item \textit{Analogy:} Walking blindfolded trying to step exactly into a small hole.
            \end{itemize}
        \end{column}
        
        \begin{column}{0.5\textwidth}
            \textbf{Optimizers (Gradient-Based)}
            \begin{itemize}
                \item E.g., SNOPT, VF13ad.
                \item Use an \textbf{Objective Function} (minimize $\Delta V$, maximize mass).
                \item Navigates constraints to find the \textit{best} possible solution, not just \textit{a} solution.
                \item \textit{Analogy:} Finding the lowest point in a valley bounded by fences.
            \end{itemize}
        \end{column}
    \end{columns}
\end{frame}

% --- Day 2 ---
\begin{frame}{Day 2: Taming Chaos}
    \framesubtitle{The Multiple Shooting Method}
    
    \textbf{The Problem with Single Shooting:}
    \begin{itemize}
        \item Propagating a single long trajectory is highly sensitive to the initial guess.
        \item A tiny $\Delta V$ error at Earth equals a massive miss distance at Mars.
    \end{itemize}
    \vspace{0.3cm}
    
    \textbf{The Multiple Shooting Architecture:}
    \begin{itemize}
        \item \textbf{Divide and Conquer:} Break the trajectory into smaller, less sensitive chunks.
        \item \textbf{Control Points:} Where the optimizer guesses the spacecraft's state.
        \item \textbf{Patch Points:} Where the algorithm enforces position and velocity continuity to "stitch" the chunks together seamlessly.
    \end{itemize}
    % Tip: Insert a graphic here showing disconnected segments merging into one curve
\end{frame}

% --- Day 3 ---
\begin{frame}{Day 3: The Low-Thrust Paradigm}
    \framesubtitle{Electric Propulsion vs. Chemical Rockets}
    
    \begin{itemize}
        \item \textbf{Breaking the Impulsive Assumption:} Standard $\Delta V$ math fails when burns take days or months instead of seconds.
        \item \textbf{Complex Orbital Changes:} 
        \begin{itemize}
            \item Changing Inclination or RAAN is difficult.
            \item Continuous burning across nodes cancels out out-of-plane changes.
            \item Requires careful, "chopped" propagation arcs.
        \end{itemize}
        \item \textbf{Why Optimization is Mandatory:} Simple targeters cannot handle the non-linear, cyclical responses of low-thrust regimes.
    \end{itemize}
\end{frame}

% --- Day 4 ---
\begin{frame}{Day 4: Optimal Control Theory}
    \framesubtitle{A Conceptual Overview}
    
    \begin{itemize}
        \item \textbf{Steering History:} We aren't solving for a single burn vector anymore; we need a continuous history of pitch and yaw over time.
        \item \textbf{The Mathematics:} Professional software (like GMAT's CSALT plugin) uses Hamiltonians and Lagrange multipliers to find these histories.
        \item \textbf{Collocation:} The solver discretizes the continuous problem into a grid of nodes to find the optimal continuously changing thrust vector.
    \end{itemize}
    \vspace{0.5cm}
    \textit{Note: The goal is to understand what the math is achieving, not necessarily to solve the differential equations by hand!}
\end{frame}

% --- Day 5 ---
\begin{frame}{Day 5: Practical Lab}
    \framesubtitle{Applying Optimization in GMAT}
    
    \textbf{Today's Mission:}
    \begin{enumerate}
        \item \textbf{The Baseline:} Set up a simple inclination change using a \textit{Targeter}. Observe how it struggles with poor initial guesses.
        \item \textbf{The Upgrade:} Swap the Targeter for an \textit{Optimizer} (e.g., VF13ad).
        \item \textbf{The Configuration:} 
        \begin{itemize}
            \item Set Objective Function: \texttt{Minimize BurnDuration}
            \item Set Constraint: Desired Inclination
        \end{itemize}
        \item \textbf{Debrief:} Analyze how the optimizer navigated the parameter space to find a fuel-efficient solution.
    \end{enumerate}
\end{frame}

% --- Conclusion ---
\begin{frame}{Summary \& Best Practices}
    \begin{itemize}
        \item \textbf{Start with "Why":} Missions like \textit{Dawn} or the \textit{Artemis Gateway} absolutely require numerical optimization; analytical math isn't enough.
        \item \textbf{The "Initial Guess" is King:} Optimizers aren't magic. They require a physically reasonable initial guess from the engineer to converge.
        \item \textbf{Software is the Standard:} Industry relies on tools like GMAT, STK, and FreeFlyer to handle these immense numerical workloads.
    \end{itemize}
\end{frame}

\end{document}